\title{Hvidløgspasta med parmesan og guanciale}
\PrepTime{25}
\CookingTime{}
\CookingTempe{}
\TypeCooking{Hurtige måltider\index[nemlig]{Hurtig!Hvidløgspasta med parmesan og guanciale}}
\NbPerson{4}
\Image{0 100 800 900}{images/nemlig/hvidloegspasta-m-parmesan-guanciale.png}

\begin{ingredient}
    \begin{main}
        \item 1 Broccoli\index[ingredient]{Broccoli}
        \item 400g Spaghetti\index[ingredient]{Spaghetti}
        \item Salt\index[ingredient]{Salt}
        \item 4 fed Hvidløg\index[ingredient]{Hvidløg}
        \item 60g Parmesan\index[ingredient]{Parmesan}
        \item 2 pk. Guanciale\index[ingredient]{Guanciale}
        \item 2 tsk. Chiliflager\index[ingredient]{Chiliflager}
        \item Peber\index[ingredient]{Peber}
    \end{main}
\end{ingredient}
\begin{recipe}
    \step{Skyl broccoli og skær i små mundrette buketter på et skærebræt (se tip!).}
    \step{Kog spaghetti i letsaltet vand i en gryde, efter pakkens anvisninger, til den er al dente.}
    \step{Gem ca. 0,5dl kogevand pr. medspiser og stil til side. Hæld resten af kogevandet fra.}
    \step{Det sidste minut tilsættes broccoli i kogevandet. Dræn vandet fra spaghetti og broccoli.}
    \step{Pil hvidløg, og snit i tynde skiver.}
    \step{Parmesan rives på et fint rivejern. Gem lidt i en lille skål til servering.}
    \step{En stor stegepande varmes op på mellem varme. På den tørre pande tilsættes guanciale, og steges langsomt til fedtet smelter.}
    \step{Når guanciale er sprødt, fiskes den af panden og stilles til side på en tallerken med køkkenrulle. Lad fedtet blive i panden.}
    \step{Tilsæt nu hvidløg og chiliflager og lad det stege til det dufter (se tip!). Det må ikke brænde på.}
    \step{Hæld spaghetti, broccoli, sprød guanciale og parmesan i panden og vend godt rundt. Krydr med salt og rigeligt peber, og justér konsistensen med lidt af det gemte kogevand.}
    \step{Server lækker italiensk spaghetti med sprødristet guanciale, parmesan og broccoli.}
\end{recipe}
\begin{notes}
    \textbf{Tip!} Du kan med fordel skylle broccolistokken, skære den i helt små tern og stege dem med på stegepanden sammen med hvidløg og chiliflager.

    \textbf{Vidste du...} At guanciale er en italiensk inspireret lufttørret kødvare lavet af svinekæbe. Den er krydret med salt og sort peber som giver en rig og intens smag til retten.
\end{notes}

\setcounter{step}{0}
\setcounter{substep}{0}

\clearpage
